\documentclass[tikz,border=4pt]{standalone}
\input{tikz-tcee}
\usepackage{pgfplotstable}
% colores-solucion.tex — Colores de los gráficos propios de la solución (1D01,
% examen de España). Se cargan después de tikz-tcee.tex (paleta «dic…»).
% Los 7 primeros son los que añade el cuadernillo del examen (mismo HEX, para
% que lo mismo tenga el mismo color en el examen y en la solución; CRITERIOS §7).
\definecolor{dicbienes}{HTML}{2C6E9B}
\definecolor{dicturismo}{HTML}{8CC3E6}
\definecolor{dicrentas}{HTML}{A6A6A6}
\definecolor{dicsaldo}{HTML}{1A1A1A}
\definecolor{dicpiin}{HTML}{8C510A}
\definecolor{dicfiscal}{HTML}{4A5A6A}
\definecolor{dicvivienda}{HTML}{B07AA1}
% Nuevos en la solución:
% sector privado (gama cálida, la del gasto privado de la guía)
\definecolor{dicprivado}{HTML}{C76A1F}
% efecto bola de nieve (azul medio, distinto del pizarra de lo fiscal)
\definecolor{dicbola}{HTML}{4F86C6}
% ajuste déficit-deuda y otras variaciones (residuos): gris claro con trama
\definecolor{dicresiduo}{HTML}{B4B4B4}
% previsión del PIB (el morado del PIB, aclarado)
\colorlet{dicpibprev}{dicpib!35!white}
\pgfplotsset{
  dicsol/.style={width=13cm, height=6.4cm, scale only axis=false,
    axis lines*=left, tick align=outside, clip=false,
    ymajorgrids, grid style={tceegris!25},
    ticklabel style={font=\footnotesize}, label style={font=\footnotesize},
    xticklabel style={/pgf/number format/1000 sep=},
    yticklabel style={/pgf/number format/.cd, use comma, 1000 sep={}},
    legend style={at={(0.5,-0.15)}, anchor=north, draw=none, font=\footnotesize,
      /tikz/every even column/.append style={column sep=8pt}},
    legend cell align=left},
  leyendalinea/.style={legend image code/.code={\draw[##1] (0cm,0cm) -- (0.55cm,0cm);}},
  % leyenda de una serie de solo puntos (en un gráfico con area legend)
  leyendamarca/.style={legend image code/.code={\draw[mark repeat=2,mark phase=2,##1] plot coordinates {(0cm,0cm) (0.3cm,0cm) (0.6cm,0cm)};}},
}

\begin{document}
% Solución, gráfico s01. Fases del ciclo de España, 2007-2027: crecimiento del
% PIB real (barras) y brecha de producción (línea); 2026-2027, previsión (barras
% claras y puntos huecos). Datos: ../datos/s01-fases.csv (cuadro 1).
% Capas (CRITERIOS §7): sombras de las fases → barras → cero → línea → puntos → rótulos.
\pgfplotstableread[col sep=comma,comment chars={\#}]{../datos/s01-fases.csv}\fases
\begin{tikzpicture}[tcee]
  \begin{axis}[dicsol, width=14.5cm, height=6.6cm,
      ybar, bar width=0.36cm,
      xmin=2006.4, xmax=2027.6, xtick={2007,...,2027},
      xticklabel style={rotate=90, anchor=east, font=\scriptsize, /pgf/number format/1000 sep={}},
      bar shift=0pt,
      ymin=-12, ymax=8, ytick={-12,-10,...,8},
      ylabel={\% y \% del PIB potencial},
      legend columns=3, legend style={at={(0.5,-0.2)}}, dicapas]
    % 1: fases (sombreados) y su nombre
    \capa{1}{
      \dicsombra[fill=tceerojo!10]{2008.5}{-12}{2013.5}{8}
      \dicsombra[fill=tceerojo!10]{2019.5}{-12}{2020.5}{8}
      \dicsombra[fill=tceegris!12]{2025.5}{-12}{2027.6}{8}
      \dicrotulo[font=\scriptsize\itshape, align=center, anchor=north]{(2007.5,7.8)}{Fin del\\auge}
      \dicrotulo[font=\scriptsize\itshape, align=center, anchor=north]{(2011,7.8)}{Doble recesión}
      \dicrotulo[font=\scriptsize\itshape, align=center, anchor=north]{(2016.5,7.8)}{Recuperación y expansión}
      \dicrotulo[font=\scriptsize\itshape, align=center, anchor=north]{(2020,7.8)}{Pan-\\demia}
      \dicrotulo[font=\scriptsize\itshape, align=center, anchor=north]{(2023,7.8)}{Rebote y expansión}
      \dicrotulo[font=\scriptsize\itshape, align=center, anchor=north]{(2026.5,7.8)}{Previsión}
    }
    % 2: crecimiento del PIB real (observado y previsto)
    \capa{2}{
      \addplot[on layer=dicbarras, dicbarra=dicpib, skip coords between index={19}{21}] table[x=anio, y=pib] {\fases};
      \addlegendentry{PIB real (variación anual, \%)}
      \addplot[on layer=dicbarras, dicbarra=dicpibprev, forget plot, skip coords between index={0}{19}] table[x=anio, y=pib] {\fases};
    }
    % 3: el cero
    \capa{3}{\dicraya{(2006.4,0) (2027.6,0)}}
    % 4: brecha de producción (línea y puntos; previsión en hueco)
    \capa{4}{
      \addplot[sharp plot, diclinea=dicsaldo, leyendalinea, mark=none] table[x=anio, y=brecha] {\fases};
      \addlegendentry{Brecha de producción (\% del PIB potencial)}
      \addplot[on layer=dicpuntos, forget plot, only marks, mark=*, mark size=1.7pt,
        mark options={fill=dicsaldo, draw=dicsaldo}, skip coords between index={19}{21}] table[x=anio, y=brecha] {\fases};
      \addplot[on layer=dicpuntos, forget plot, only marks, mark=*, mark size=1.7pt,
        mark options={fill=white, draw=dicsaldo, line width=0.8pt}, skip coords between index={0}{19}] table[x=anio, y=brecha] {\fases};
    }
    % 5: los datos que cuentan
    \capa{5}{
      \dicrotulo[font=\scriptsize, anchor=west, inner sep=1pt]{(2007.15,4.9)}{brecha +4,9}
      \dicrotulo[font=\scriptsize, anchor=north, inner sep=1pt]{(2013,-7.1)}{brecha $-$6,9}
      \dicrotulo[font=\scriptsize, anchor=south east, inner sep=1pt]{(2019,2.8)}{brecha +2,6}
      \dicrotulo[font=\scriptsize, anchor=north, inner sep=1pt]{(2020,-11.0)}{PIB $-$10,9}
    }
  \end{axis}
\end{tikzpicture}
\end{document}
